\chapter{Onsager--Kaufman 精确谱、有限扇区与自由能}
\label{chap:phys-sm-onsager-exact-solution}

Jordan--Wigner 与 Clifford 线性化已经把 transfer operator 化成独立动量块，同时保留了 NS/R 网格和物理奇偶投影。从这些有限尺寸对象出发，先得到单粒子 rapidity \(\gamma(q)\)，再写出四个 fermionic traces 的有限乘积，最后取热力学极限得到 Onsager 自由能。临界线由单粒子 transfer mass 闭合确定；具体自旋相关长度还需结合观测算符的 form-factor selection rule，不能与单粒子谱隙直接等同。

\section{动量块的 Clifford 线性变换}
\label{sec:phys-sm-onsager-solution}

在固定 NS 或 R 空间扇区内，Fourier 变换把平移不变二次型分成 \(\{q,-q\}\) 块。选取一个 Nambu/Majorana 基后，垂直和水平 transfer factors 的共轭作用分别由 determinant-one 矩阵
\begin{equation}
 \mathsf R_v=
 \begin{pmatrix}
 \cosh 2K_v^*&\sinh 2K_v^*\\
 \sinh 2K_v^*&\cosh 2K_v^*
 \end{pmatrix}
 \label{eq:phys-sm-rv-block}
\end{equation}
以及
\begin{equation}
 \mathsf R_h(q)=
 \begin{pmatrix}
 \cosh 2K_h&-\ee^{-\ii q}\sinh 2K_h\\
 -\ee^{\ii q}\sinh 2K_h&\cosh 2K_h
 \end{pmatrix}
 \label{eq:phys-sm-rh-block}
\end{equation}
表示。一个完整对称 transfer step 与 \(\mathsf R_v\mathsf R_h(q)\) 相似，因此具有相同本征值。令
\[
 \mathsf R(q)=\mathsf R_v\mathsf R_h(q).
\]
因为 \(\det\mathsf R(q)=1\)，本征值可写为 \(\ee^{\pm\gamma(q)}\)。半迹给出
\begin{equation}
 \boxed{
 \cosh\gamma(q)
 =\cosh(2K_v^*)\cosh(2K_h)
 -\sinh(2K_v^*)\sinh(2K_h)\cos q.
 }
 \label{eq:phys-sm-gamma-dispersion}
\end{equation}
取 \(\gamma(q)\ge0\) 定义物理解支。

\begin{derivation}[从 Clifford 生成元到 \(2\times2\) Nambu 块]
固定 NS 或 R 扇区后引入 Majorana 算符 \(\xi_{2j-1},\xi_{2j}\)，满足 \(\{\xi_a,\xi_b\}=2\delta_{ab}\)。Jordan--Wigner 映射把 \(\tau_j^z\) 写为 \(\xi_{2j-1}\xi_{2j}\)，把 \(\tau_j^x\tau_{j+1}^x\) 写为 \(\xi_{2j}\xi_{2j+1}\)，因此垂直生成元 \(K_v^*\sum_j\tau_j^z\) 在 Majorana 变量中是逐站点双线性 \(K_v^*\sum_j\xi_{2j-1}\xi_{2j}\)，水平生成元 \(K_h\sum_j\tau_j^x\tau_{j+1}^x\) 是最近邻双线性 \(K_h\sum_j\xi_{2j}\xi_{2j+1}\)。

由\eqnrefs{eq:phys-sm-clifford-linear-action}，指数共轭作用 \(\ee^{-G}\xi\ee^{G}\) 对双线性 \(G=K\sum_j\xi_{2j+a}\xi_{2j+b}\) 化为各 \((\xi_{2j+a},\xi_{2j+b})\) 二维子空间上的双曲旋转。对垂直项，每个站点独立旋转角 \(2K_v^*\)，故
\begin{equation*}
 \xi_{2j-1}\mapsto \cosh(2K_v^*)\xi_{2j-1}+\sinh(2K_v^*)\xi_{2j},
 \qquad
 \xi_{2j}\mapsto \sinh(2K_v^*)\xi_{2j-1}+\cosh(2K_v^*)\xi_{2j},
\end{equation*}
矩阵即\eqnrefs{eq:phys-sm-rv-block}。对水平项，最近邻耦合链接 \(j\) 与 \(j+1\)，先作 Fourier 变换 \(\xi_{2j}=\frac1{\sqrt N}\sum_q\ee^{\ii qj}\tilde\xi_2(q)\)，于是
\begin{align*}
 \sum_j\xi_{2j}\xi_{2j+1}
 &=\frac1N\sum_{j,q,q'}\ee^{\ii qj}\ee^{\ii q'(j+1)}\tilde\xi_2(q)\tilde\xi_2(q')\\
 &=\sum_q\ee^{\ii q}\tilde\xi_2(q)\tilde\xi_2(-q)
 =\sum_q\frac{\ee^{\ii q}+\ee^{-\ii q}}2\,\tilde\xi_2(q)\tilde\xi_2(-q),
\end{align*}
其中第二步用了 \(\frac1N\sum_j\ee^{\ii(q+q')j}=\delta_{q,-q'}\)，第三步把 \(q\mapsto-q\) 两份合并。在每个 \(\{q,-q\}\) 二维不变子空间中，\(\ee^{\pm\ii q}\) 相位恰给出\eqnrefs{eq:phys-sm-rh-block}。

把 \(\mathsf R_v\) 与 \(\mathsf R_h(q)\) 直接相乘，
\begin{align*}
 \mathsf R(q)
 &=
 \begin{pmatrix}
 \cosh 2K_v^*&\sinh 2K_v^*\\
 \sinh 2K_v^*&\cosh 2K_v^*
 \end{pmatrix}
 \begin{pmatrix}
 \cosh 2K_h&-\ee^{-\ii q}\sinh 2K_h\\
 -\ee^{\ii q}\sinh 2K_h&\cosh 2K_h
 \end{pmatrix}\\
 &=
 \begin{pmatrix}
 \cosh 2K_v^*\cosh 2K_h-\ee^{\ii q}\sinh 2K_v^*\sinh 2K_h
 &-\ee^{-\ii q}\cosh 2K_v^*\sinh 2K_h+\sinh 2K_v^*\cosh 2K_h\\
 \sinh 2K_v^*\cosh 2K_h-\ee^{\ii q}\cosh 2K_v^*\sinh 2K_h
 &-\ee^{-\ii q}\sinh 2K_v^*\sinh 2K_h+\cosh 2K_v^*\cosh 2K_h
 \end{pmatrix}.
\end{align*}
对角元相加并使用 \(\ee^{\ii q}+\ee^{-\ii q}=2\cos q\)，
\begin{align*}
 \Tr\mathsf R(q)
 &=2\cosh(2K_v^*)\cosh(2K_h)
 -\sinh(2K_v^*)\sinh(2K_h)(\ee^{\ii q}+\ee^{-\ii q})\\
 &=2\left[
 \cosh(2K_v^*)\cosh(2K_h)
 -\sinh(2K_v^*)\sinh(2K_h)\cos q
 \right].
\end{align*}
又因 \(\cosh^2-\sinh^2=1\)，\(\det\mathsf R_v=\det\mathsf R_h(q)=1\)，故 \(\det\mathsf R(q)=1\)。设两本征值为 \(\lambda,\lambda^{-1}\)，于是 \(\Tr\mathsf R=\lambda+\lambda^{-1}\)。令 \(\lambda=\ee^{\gamma}\)，则 \(\lambda+\lambda^{-1}=\ee^{\gamma}+\ee^{-\gamma}=2\cosh\gamma\)。对照上式即得\eqnrefs{eq:phys-sm-gamma-dispersion}。整个推导只用 Majorana 反对易关系、Fourier 变换与 \(2\times2\) 矩阵乘法，未涉及热力学极限。
\end{derivation}

\section{有限尺寸谱与四个精确 traces}

定义单步整体标量
\[
 \mathcal C_N=[2\sinh(2K_v)]^{N/2}.
\]
在固定空间扇区 \(s\in\{\mathrm{NS},\mathrm R\}\) 中，二次 transfer operator 可作 Bogoliubov 正规形，使多体本征值写成
\begin{equation}
 \Lambda_s(\{n_q\})
 =\mathcal C_N
 \exp\!\left[
 \frac12\sum_{q\in s}(1-2n_q)\gamma(q)
 \right],
 \qquad n_q\in\{0,1\}.
 \label{eq:phys-sm-finite-spectrum-structure}
\end{equation}
这里对自反模也把 \(n_q\) 视为独立单模占据；非自反 \(q,-q\) 的 canonical transformation 则成对组织。令 \(p_s=\pm1\) 表示 Bogoliubov 真空相对于原 Jordan--Wigner 奇偶算符的奇偶，即
\[
 \mathcal P
 =p_s(-1)^{\sum_q n_q}.
\]
非自反的 \(\{q,-q\}\) Bogoliubov 变换只成对改变原费米子占据，因此不会改变真空奇偶。自反模中，\(q=\pi\) 的未配对质量具有符号 \(K_v^*+K_h>0\)，而 R 扇区的 \(q=0\) 未配对质量具有符号 \(K_v^*-K_h\)。于是对铁磁耦合 \(K_h,K_v>0\) 且离开临界线，当前 Jordan--Wigner convention 给出
\begin{equation}
 p_{\rm NS}=+1,
 \qquad
 p_{\rm R}=\operatorname{sgn}(K_v^*-K_h).
 \label{eq:phys-sm-bogoliubov-vacuum-parity}
\end{equation}
在临界线 \(K_v^*=K_h\) 上，R 扇区的 \(q=0\) rapidity 为零，此时含奇偶插入的 R trace 本身为零，所以 \(p_{\rm R}\) 在该点无需单独指定。

对纵向周期数 \(M\)，四个 trace 由独立 quasiparticle 模逐项相乘：
\begin{align}
 \mathcal Z_{s,0}
 &=\mathcal C_N^M
 \prod_{q\in s}
 2\cosh\frac{M\gamma(q)}{2},
 \label{eq:phys-sm-sector-trace-cosh}\\
 \mathcal Z_{s,1}
 &=p_s\mathcal C_N^M
 \prod_{q\in s}
 2\sinh\frac{M\gamma(q)}{2}.
 \label{eq:phys-sm-sector-trace-sinh}
\end{align}
物理环面配分函数再由奇偶投影\eqnrefs{eq:phys-sm-kaufman-sector-projection}
\begin{equation}
 Z_{N,M}
 =\frac12(\mathcal Z_{\mathrm{NS},0}+\mathcal Z_{\mathrm{NS},1})
 +\frac12(\mathcal Z_{\mathrm R,0}-\mathcal Z_{\mathrm R,1})
 \label{eq:phys-sm-finite-torus-four-traces}
\end{equation}
得到。若存在零 rapidity，自带 \(\sinh(M\gamma/2)=0\) 会自动使相应奇偶插入 trace 消失；这正是不能把特殊模藏在“所有 \(q\) 成对”近似中的原因。

\begin{derivation}[为什么奇偶插入把 \(\cosh\) 变成 \(\sinh\)]
由\eqnrefs{eq:phys-sm-finite-spectrum-structure}，固定扇区 \(s\) 中单步 transfer 的多体本征值为
\begin{equation*}
 \Lambda_s(\{n_q\})
 =\mathcal C_N\exp\!\left[\frac12\sum_{q\in s}(1-2n_q)\gamma(q)\right],
 \qquad n_q\in\{0,1\},
\end{equation*}
其中 \(n_q=0\) 对应 Bogoliubov 真空，\(n_q=1\) 对应单模激发。纵向 \(M\) 步等价于对每个 quasiparticle 模作 \(M\) 次乘幂，故对单模 \(q\)，
\begin{equation*}
 \Lambda_q(n)^M
 =\mathcal C_N^M\exp\!\left[\frac{M\gamma(q)}2(1-2n)\right].
\end{equation*}
对 \(n=0,1\) 求和，未插入奇偶时
\begin{align*}
 \sum_{n=0}^1\Lambda_q(n)^M
 &=\mathcal C_N^M\left[
 \exp\!\left(+\frac{M\gamma(q)}2\right)
 +\exp\!\left(-\frac{M\gamma(q)}2\right)
 \right]\\
 &=\mathcal C_N^M\,2\cosh\frac{M\gamma(q)}2,
\end{align*}
第二步用了 \(\cosh x=(\ee^x+\ee^{-x})/2\)。若迹中插入原费米子奇偶 \((-1)^n\)，则
\begin{align*}
 \sum_{n=0}^1(-1)^n\Lambda_q(n)^M
 &=\mathcal C_N^M\left[
 \exp\!\left(+\frac{M\gamma(q)}2\right)
 -\exp\!\left(-\frac{M\gamma(q)}2\right)
 \right]\\
 &=\mathcal C_N^M\,2\sinh\frac{M\gamma(q)}2,
\end{align*}
第二步用了 \(\sinh x=(\ee^x-\ee^{-x})/2\)。全部模独立相乘，得到无奇偶插入与有奇偶插入的扇区迹
\begin{align*}
 \mathcal Z_{s,0}
 &=\mathcal C_N^M\prod_{q\in s}2\cosh\frac{M\gamma(q)}2,\\
 \mathcal Z_{s,1}
 &=p_s\,\mathcal C_N^M\prod_{q\in s}2\sinh\frac{M\gamma(q)}2,
\end{align*}
即\eqnrefs{eq:phys-sm-sector-trace-cosh}--\eqnrefs{eq:phys-sm-sector-trace-sinh}。其中 \(p_s\) 是 Bogoliubov 真空相对于原 Jordan--Wigner 奇偶算符的符号：原费米子奇偶为
\begin{equation*}
 \mathcal P=p_s(-1)^{\sum_q n_q},
\end{equation*}
因此插入 \(\mathcal P\) 后每个模贡献 \((-1)^{n_q}\)，整体再乘 \(p_s\)。

对非自反块 \(\{q,-q\}\)，canonical transformation 混合的是 \(q\) 与 \(-q\) 两个模，真空占据改变为偶数，故不翻转 \(p_s\)；只有自反模能够翻转。自反模 \(q=\pi\) 的未配对质量符号为 \(K_v^*+K_h>0\)，不穿零；\(q=0\) 的符号为 \(K_v^*-K_h\)，恰在临界线 \(K_v^*=K_h\) 穿零。因此铁磁耦合下
\begin{equation*}
 p_{\rm NS}=+1,
 \qquad
 p_{\rm R}=\operatorname{sgn}(K_v^*-K_h),
\end{equation*}
即\eqnrefs{eq:phys-sm-bogoliubov-vacuum-parity}。整个论证在有限 \(M,N\) 上成立，不需要热力学极限。
\end{derivation}

对长柱面自由能只需各扇区主导本征值的体积阶。奇偶投影最多改变有限个 quasiparticle 的占据，所以
\begin{equation}
 \log\Lambda_{\max}(N)
 =\frac N2\log[2\sinh(2K_v)]
 +\frac12\sum_{q\in s}\gamma(q)+o(N),
 \label{eq:phys-sm-lambda-max-finite}
\end{equation}
其中任一物理主导网格在 \(N\to\infty\) 都给出同一 Riemann 积分。

\section{Onsager 自由能、临界线与 transfer mass}

将\eqnrefs{eq:phys-sm-lambda-max-finite} 除以 \(N\) 并取 \(N\to\infty\)，得到各向异性零场方格 Ising 的单积分形式
\begin{equation}
 \boxed{
 -\beta f
 =\frac12\log[2\sinh(2K_v)]
 +\frac{1}{2\pi}\int_0^\pi\gamma(q)\,\dd q.
 }
 \label{eq:phys-sm-onsager-standard}
\end{equation}
利用 \(\sinh(2K_v)\sinh(2K_v^*)=1\) 可代数变形为常见的对称双积分表示；二者描述同一自由能。

色散最小值位于 \(q=0\)。因为
\[
 \cosh\gamma(0)
 =\cosh[2(K_v^*-K_h)],
\]
所以
\begin{equation}
 \Delta_{\rm tr}:=\gamma(0)=2|K_v^*-K_h|.
 \label{eq:phys-sm-gapclose}
\end{equation}
这里 \(\Delta_{\rm tr}\) 是 dimensionless elementary transfer mass。闭隙条件 \(K_v^*=K_h\) 与
\begin{equation}
 \boxed{
 \sinh(2K_h)\sinh(2K_v)=1
 }
 \label{eq:phys-sm-critical-line}
\end{equation}
等价，从而独立验证 Kramers--Wannier 自对偶候选线确实是谱隙闭合线。

临界点附近令 \(\Delta_K=2(K_v^*-K_h)\)。由\eqnrefs{eq:phys-sm-gamma-dispersion} 展开，
\begin{equation}
 \gamma(q)^2
 =\Delta_K^2+v_c^2q^2
 +O(\Delta_K^4,\Delta_K^2q^2,q^4),
 \qquad
 v_c=\sinh(2K_{h,c})>0.
 \label{eq:phys-sm-critical-dispersion}
\end{equation}
若原始耦合 \(J_h,J_v\) 不随温度改变，则
\begin{equation}
 \Delta_K=a_T(T-T_c)+O((T-T_c)^2),
 \qquad
 a_T=\frac{2}{T_c}
 \left[
 K_{h,c}+\frac{K_{v,c}}{\sinh(2K_{v,c})}
 \right]>0.
 \label{eq:phys-sm-gap-temperature-slope}
\end{equation}
所以 \(\Delta_{\rm tr}\) 线性闭合。

\begin{derivation}[临界自由能的对数奇性]
自由能的非解析部分只来自低 $q$ 模式的积分。在临界点附近，色散关系 $\omega(q)\approx\sqrt{\Delta_K^2+v_c^2q^2}$ 产生 $\sqrt{\Delta_K^2+v_c^2q^2}-v_cq$ 的被积函数，其对 $\Delta_K^2$ 的导数含 $\operatorname{arsinh}$，积分后产生 $\Delta_K^2\log|\Delta_K|$ 型奇性。

局部模型。减去临界点的光滑背景，自由能非解析部分只来自小 $q$ 区域：
\begin{equation*}
 \mathcal J(\Delta_K)
 =\int_0^{q_0}
 \left(\sqrt{\Delta_K^2+v_c^2q^2}-v_cq\right)\dd q.
\end{equation*}
被积函数中 $-v_cq$ 扣除临界点 $\Delta_K=0$ 时的光滑背景 $\sqrt{v_c^2q^2}=v_cq$，剩余部分在 $q\to0$、$\Delta_K\to0$ 时产生奇性。

对 $\Delta_K^2$ 求导。
\begin{equation*}
 \frac{\partial\mathcal J}{\partial(\Delta_K^2)}
 =\int_0^{q_0}\frac{1}{2\sqrt{\Delta_K^2+v_c^2q^2}}\dd q
 =\frac{1}{2v_c}\operatorname{arsinh}\frac{v_cq_0}{|\Delta_K|}.
\end{equation*}
当 $|\Delta_K|\to0$，$\operatorname{arsinh}(v_cq_0/|\Delta_K|)\sim\log(2v_cq_0/|\Delta_K|)$，故
\begin{equation*}
 \frac{\partial\mathcal J}{\partial(\Delta_K^2)}
 =-\frac{1}{2v_c}\log|\Delta_K|+O(1).
\end{equation*}

积分回去。对 $\Delta_K^2$ 积分，
\begin{equation*}
 \mathcal J(\Delta_K)
 =-\frac{\Delta_K^2}{2v_c}\log|\Delta_K|
 +O(\Delta_K^2).
\end{equation*}
自由能奇异部是 $\Delta_K^2\log|\Delta_K|$ 型。

温度依赖。由 $K_h=J_h/(\kB T)$、$K_v=J_v/(\kB T)$，对偶关系 $\sinh(2K_v)\sinh(2K_v^*)=1$ 微分给出
\begin{equation*}
 \frac{\dd K_v^*}{\dd K_v}
 =-\frac{1}{\sinh(2K_v)}.
\end{equation*}
因此 $\Delta_K$ 对 $T$ 的导数在 $T_c$ 处为
\begin{equation*}
 a_T=\frac{2}{T_c}
 \left[K_{h,c}+\frac{K_{v,c}}{\sinh(2K_{v,c})}\right]>0,
\end{equation*}
即\eqnrefs{eq:phys-sm-gap-temperature-slope}，故 $\Delta_K=a_T(T-T_c)+O((T-T_c)^2)$。

自由能奇异部是 $(T-T_c)^2\log|T-T_c|$，二阶温度导数表现为 $\log|T-T_c|$ 型对数发散——这正是二维 Ising 模型的著名对数比热奇性。
\end{derivation}

自由能奇性由最软 transfer mode 控制，但具体观测量的相关长度还取决于它能耦合到哪些激发。因而在进入有限尺寸临界修正前，需要把单粒子 transfer mass 与自旋相关长度严格区分。

由 transfer 谱得到的 \(\Delta_{\rm tr}=\gamma(0)\) 首先只是最低 quasiparticle mass。具体观测量的相关长度还要使用\eqnrefs{eq:phys-sm-observable-correlation-length} 的 selection rule。二维 Ising 的精确 form-factor theory 给出额外输入：无序相中的 spin field 可耦合到一粒子通道，而有序纯相的 connected spin correlation 首个非真空贡献是两粒子通道。因此在接近临界点时
\begin{equation}
 \xi_{+,\sigma}^{-1}=\Delta_{\rm tr},
 \qquad
 \xi_{-,\sigma,c}^{-1}=2\Delta_{\rm tr},
 \label{eq:phys-sm-spin-correlation-masses}
\end{equation}
这里第二式指有序纯相的 connected correlation。非零 form factor 本身是独立的 Ising correlation-theory 输入；仅凭自由费米单粒子谱不能证明其系数非零。两侧相关长度振幅可以不同，而 \(\Delta_{\rm tr}=|a_T|\,|T-T_c|+O((T-T_c)^2)\) 仍给出共同临界指数 \(\nu=1\)。

\section{临界 NS/R 和的有限尺寸修正}

有限扇区在临界点不会只留下“\(O(1/N)\)”的模糊量级；系数可以直接从各向异性色散的 cusp 算出。在临界线 \(K_v^*=K_h\) 上，令
\[
 v_c:=\sinh(2K_{h,c})>0.
\]
由\eqnrefs{eq:phys-sm-gamma-dispersion}
\[
 \cosh\gamma_c(q)
 =1+2v_c^2\sin^2\frac q2,
\]
所以有精确临界 rapidity
\begin{equation}
 \boxed{
 \gamma_c(q)=2\operatorname{arsinh}
 \left(v_c\left|\sin\frac q2\right|\right).
 }
 \label{eq:phys-sm-critical-anisotropic-rapidity}
\end{equation}
各向同性临界点只是 \(v_c=1\) 的特例。定义
\[
 S_{\rm R}(N)=\sum_{m=0}^{N-1}
 \gamma_c\!\left(\frac{2\pi m}{N}\right),
 \qquad
 S_{\rm NS}(N)=\sum_{m=0}^{N-1}
 \gamma_c\!\left(\frac{(2m+1)\pi}{N}\right),
\]
以及平均值 \(\bar\gamma=(2\pi)^{-1}\int_0^{2\pi}\gamma_c(q)\dd q\)。则
\begin{align}
 S_{\rm R}(N)
 &=N\bar\gamma-\frac{\pi v_c}{3N}+O(N^{-3}),
 \label{eq:phys-sm-critical-r-sum}\\
 S_{\rm NS}(N)
 &=N\bar\gamma+\frac{\pi v_c}{6N}+O(N^{-3}).
 \label{eq:phys-sm-critical-ns-sum}
\end{align}
所以
\begin{equation}
 S_{\rm NS}(N)-S_{\rm R}(N)
 =\frac{\pi v_c}{2N}+O(N^{-3}).
 \label{eq:phys-sm-critical-sector-splitting}
\end{equation}
主导 NS 条带因此满足
\begin{equation}
 -\beta f_N
 =-\beta f_\infty
 +\frac{\pi v_c}{12N^2}+O(N^{-4}).
 \label{eq:phys-sm-critical-free-energy-fss}
\end{equation}
各向异性只通过临界 transfer velocity \(v_c\) 改变度量尺度；无量纲系数仍与周期 Ising 临界理论的 central charge \(c=1/2\) 一致。这里并未把共形场论作为推导输入：系数直接来自有限格点 rapidity 的 cusp 与离散 Fourier 取样。

\begin{derivation}[各向异性临界 cusp 到 NS/R 的 \(1/N\) 系数]
把 \(\gamma_c\) 作 \(2\pi\)-周期 Fourier 展开
\[
 \gamma_c(q)=\sum_{n\in\mathbb Z}
 \widehat\gamma_n\ee^{\ii nq}.
\]
由\eqnrefs{eq:phys-sm-critical-anisotropic-rapidity}，在 \(q=0\pmod{2\pi}\) 附近
\[
 \gamma_c(q)=v_c|q|+O(|q|^3).
\]
因此一阶导数的周期跳跃为 \(2v_c\)。分部积分两次，连续性使第一次的边界项消失，第二次由导数跳跃给出
\begin{equation}
 \widehat\gamma_n
 =-\frac{v_c}{\pi n^2}+O(n^{-4}),
 \qquad |n|\to\infty.
 \label{eq:phys-sm-critical-gamma-fourier-tail}
\end{equation}
离散 Fourier 取样恒等式为
\[
 S_{\rm R}=N\sum_{\ell\in\mathbb Z}
 \widehat\gamma_{\ell N},
 \qquad
 S_{\rm NS}=N\sum_{\ell\in\mathbb Z}
 (-1)^\ell\widehat\gamma_{\ell N}.
\]
去掉 \(\ell=0\) 的 bulk 项并使用
\[
 \sum_{\ell=1}^\infty\frac1{\ell^2}=\frac{\pi^2}{6},
 \qquad
 \sum_{\ell=1}^\infty\frac{(-1)^\ell}{\ell^2}=-\frac{\pi^2}{12},
\]
得到
\[
 \begin{aligned}
 S_{\rm R}-N\bar\gamma
 &=-\frac{2v_c}{\pi N}
 \sum_{\ell=1}^\infty\frac1{\ell^2}
 +O(N^{-3})
 =-\frac{\pi v_c}{3N}+O(N^{-3}),\\
 S_{\rm NS}-N\bar\gamma
 &=-\frac{2v_c}{\pi N}
 \sum_{\ell=1}^\infty\frac{(-1)^\ell}{\ell^2}
 +O(N^{-3})
 =\frac{\pi v_c}{6N}+O(N^{-3}).
 \end{aligned}
\]
这证明\eqnrefs{eq:phys-sm-critical-r-sum}--\eqnrefs{eq:phys-sm-critical-sector-splitting}。由于 \(\log\Lambda\) 含 \(S_s/2\)，再除以宽度 \(N\) 就得到\eqnrefs{eq:phys-sm-critical-free-energy-fss} 的 \(N^{-2}\) 修正。
\end{derivation}

传递矩阵的 Clifford 谱、有限 Fock 空间中的迹、热力学自由能和临界有限尺寸求和均由前述有限尺寸算符直接推出。自旋二点函数本身不是二次费米子局域算符；把它化为 Toeplitz 行列式或 Pfaffian 仍需独立的相关函数定理。结合这一输入与 strong Szegő 定理后，Yang 自发磁化可由剩余 Fourier 求和得到。

\begin{derivation}[Toeplitz 行列式与 Montroll--Ward--Kaufman 关联函数]
自旋二点函数 \(\langle\sigma_{0,0}\sigma_{M,N}\rangle\) 可化为 Toeplitz 行列式，其渐近由 Szegő 定理给出。分四步建立。

Kaufman 旋转算符表示。Kaufman 把 Ising 自旋算符表示为 Clifford 代数中的投影算符。对一行内 \(N\) 个自旋，定义旋转算符
\begin{equation*}
  \mathbf W=\prod_{j=1}^N\frac{1+\tau_j^x}{\sqrt2},
\end{equation*}
使 \(\sigma_j^z\mathbf W=\mathbf W\tau_j^x\)（把 \(Z\) 旋转到 \(X\) 方向）。两点函数化为 transfer matrix 的矩阵元
\begin{equation*}
  \langle\sigma_{0,0}\sigma_{M,N}\rangle=\frac{\langle0|\mathbf W^\dagger\mathsf T^M\boldsymbol\sigma_N\mathsf T^M\mathbf W|0\rangle}{\langle0|\mathsf T^{2M}|0\rangle},
\end{equation*}
其中 \(\boldsymbol\sigma_N=\prod_{j=1}^N\tau_j^x\) 是一列自旋算符的 Pauli 表示。

Clifford 表示与 Toeplitz 结构。Jordan--Wigner 变换后 \(\boldsymbol\sigma_N=\ii^{N/2}\prod_{j=1}^N(c_j+c_j^\dagger)\)（弦算符乘积）。在自由费米子 Fock 空间中，这一乘积的期望值化为 \(N\times N\) 行列式，其元素是费米子 Green 函数
\begin{equation*}
  G_{jk}=\langle0|c_j\mathsf T^M c_k^\dagger\mathsf T^M|0\rangle/\Lambda_0^{2M}.
\end{equation*}
平移不变性使 \(G_{jk}=G_{j-k}\)（只依赖差），行列式成为 Toeplitz 行列式 \(D_N=\det[G_{j-k}]_{j,k=1}^N\)。Green 函数的 Fourier 表示
\begin{equation*}
  G_n=\frac{1}{2\pi}\int_0^{2\pi}\Phi(\theta)\ee^{-\in\theta}\dd\theta
\end{equation*}
定义 Toeplitz 符号 \(\Phi(\theta)\)，它由 Ising 耦合和距离 \((M,N)\) 通过 BdG 谱决定。

strong Szegő 定理。对充分光滑且零环绕数（winding number \(=0\)）的符号 \(\Phi\)，strong Szegő 定理给出
\begin{equation}
  \log D_N=N\log G_0+\sum_{n=1}^\infty n(\log\Phi)_n(\log\Phi)_{-n}+o(1),
  \label{eq:phys-sm-szego-strong}
\end{equation}
其中 \((\log\Phi)_n\) 是 \(\log\Phi(\theta)\) 的 Fourier 系数。第一项是"体积项"（正比于 \(N\)），第二项是 Szegő "常数项"（不依赖 \(N\)）。对 Ising 符号函数，\(\log\Phi\) 的 Fourier 系数可由 BdG 色散 \(\gamma(q)\) 显式计算。

Yang 磁化公式的推导。自发磁化 \(m=\lim_{N\to\infty}\langle\sigma_{0,0}\sigma_{M,N}\rangle^{1/(MN)}\)（取适当极限）。由\eqnrefs{eq:phys-sm-szego-strong}，体积项给出 \(\log m=\log G_0\)，常数项给出有限尺寸修正。对 Ising 符号 \(\Phi(\theta)=[1-k^2\sin^2(\theta/2)]^{1/2}\)（\(k=\sinh^{-2}(2K)\)），Fourier 系数
\begin{equation*}
  (\log\Phi)_n=-\frac{1}{2n}\left(\frac{1-\sqrt{1-k^2}}{k}\right)^{|n|},
\end{equation*}
求和 \(\sum_{n=1}^\infty n(\log\Phi)_n(\log\Phi)_{-n}=\frac{1}{4}\log(1-k^2)\)，给出\eqnrefs{eq:phys-sm-onsager-magnetization}。\(T>T_c\) 时 \(k>1\)，符号 \(\Phi\) 有零点（winding number \(\neq0\)），strong Szegő 定理失效，需用 generalized Szegő 定理（Fisher--Hartwig），给出磁化为零。

临界点 \(K=K_c\)（\(k=1\)）时 \(\Phi(\theta)\) 在 \(\theta=0\) 有 cusp 奇点，strong Szegő 定理不适用。Fisher--Hartwig 把符号分解为 \(\Phi=\Phi_B\prod_j t_{\alpha_j,\beta_j}\)，其中 \(\Phi_B\) 满足 Szegő 条件，\(t_{\alpha,\beta}\) 是奇点因子。Ising 临界对应单个 \((\alpha,\beta)=(0,1/2)\) 奇点，行列式渐近
\begin{align*}
  D_N
  &\sim E\,N^{\alpha^2-\beta^2}
  \\
  &\quad{}\times\exp\!\left[
    N\log G_0+\sum n(\log\Phi_B)_n(\log\Phi_B)_{-n}
  \right],
\end{align*}
幂律 \(N^{-1/4}\) 给出 \(\eta=1/4\)。这把 Szegő/Fisher--Hartwig 渐近与 CFT 标度维度 \(h_\sigma=1/16\)（因 \(2h_\sigma=1/8\)，\(\eta=2h_\sigma=1/4\)）精确对应。
\end{derivation}
